\section{Temporal Discretization}
\label{sec:temporal_discretization}

\subsection{Approximate Factorization}
\label{subsec:approximate_factorization}

\subsubsection{Implicit Euler Discretization}

Following the spatial discretization presented in the previous
section, the energy equation can be expressed as a system of
ordinary differential equations:
\begin{equation}
\frac{dT_{ij}}{dt}=FI_{ij}+S_{ij}=RHS_{ij},
\label{eq:semi_discrete_energy}
\end{equation}
where $FI_{ij}$ denotes the combined advective and diffusive
contributions, and $S_{ij}$ is the viscous-dissipation source
term. The cell-centered temperature is denoted by $T_{ij}$,
with the control-volume average represented by the same quantity
under the present second-order approximation.

The implicit Euler method approximates the temporal derivative
using a backward difference and evaluates the spatial operator
at the next time level:
\begin{equation}
\frac{T_{ij}^{n+1}-T_{ij}^{n}}{\Delta t}
=
FI_{ij}^{n+1}+S_{ij}^{n+1}.
\label{eq:implicit_euler}
\end{equation}
Here, $n$ and $n+1$ denote the current and next time levels,
respectively, and $\Delta t$ is the time step. Since the
velocity field is prescribed in the present problem, the
viscous-dissipation source term is independent of temperature
and can be evaluated directly from the velocity gradients.

Substituting the spatially discretized expression into
Equation~\eqref{eq:implicit_euler} gives
\begin{equation}
\begin{aligned}
\frac{T_{ij}^{n+1}-T_{ij}^{n}}{\Delta t}
={}&
-\frac{
u_{i+1,j}T_{i+1,j}^{n+1}
-u_{i-1,j}T_{i-1,j}^{n+1}
}{2\Delta x}
\\[0.15cm]
&+
\frac{1}{\mathrm{Re}\,\mathrm{Pr}}
\frac{
T_{i+1,j}^{n+1}-2T_{ij}^{n+1}
+T_{i-1,j}^{n+1}
}{\Delta x^2}
\\[0.15cm]
&-
\frac{
v_{i,j+1}T_{i,j+1}^{n+1}
-v_{i,j-1}T_{i,j-1}^{n+1}
}{2\Delta y}
\\[0.15cm]
&+
\frac{1}{\mathrm{Re}\,\mathrm{Pr}}
\frac{
T_{i,j+1}^{n+1}-2T_{ij}^{n+1}
+T_{i,j-1}^{n+1}
}{\Delta y^2}
+S_{ij}.
\end{aligned}
\label{eq:implicit_euler_expanded}
\end{equation}

Although the implicit Euler method provides a direct temporal
discretization, Equation~\eqref{eq:implicit_euler_expanded}
couples the unknown temperatures at the next time level.
Solving the resulting system simultaneously for all control
volumes can be computationally expensive. The approximate
factorization method addresses this difficulty by splitting
the two-dimensional operator into directional operators,
allowing the solution to be obtained through successive
one-dimensional solves.

\subsubsection{Delta Formulation}

To formulate the approximate factorization method, the
temperature update between successive time levels is defined as
\begin{equation}
\delta T_{ij}^{n+1}
=
T_{ij}^{n+1}-T_{ij}^{n}.
\label{eq:temperature_update}
\end{equation}
Thus, the temperature at the next time level can be written as
\begin{equation}
T_{ij}^{n+1}=T_{ij}^{n}+\delta T_{ij}^{n+1}.
\label{eq:temperature_update_relation}
\end{equation}

Substituting Equation~\eqref{eq:temperature_update_relation}
into the implicit Euler equation and separating the terms
involving the unknown update from those evaluated at the
current time level yields
\begin{equation}
\begin{aligned}
\frac{\delta T_{ij}^{n+1}}{\Delta t}
+\left(\mathcal{X}\delta T^{n+1}\right)_{ij}
+\left(\mathcal{Y}\delta T^{n+1}\right)_{ij}
=
RHS_{ij}^{n},
\end{aligned}
\label{eq:delta_form}
\end{equation}
where
\begin{equation}
RHS_{ij}^{n}=FI_{ij}^{n}+S_{ij}.
\label{eq:rhs_old_time}
\end{equation}
Here, $FI_{ij}^{n}$ is evaluated using the temperature field
at time level $n$, and $S_{ij}$ is the prescribed
viscous-dissipation source term.

The directional operators $\mathcal{X}$ and $\mathcal{Y}$
contain the advective and diffusive contributions in the
$x$- and $y$-directions, respectively. They are defined by
\begin{equation}
\begin{aligned}
\left(\mathcal{X}T\right)_{ij}
={}&
\frac{
u_{i+1,j}T_{i+1,j}
-u_{i-1,j}T_{i-1,j}
}{2\Delta x}
\\
&-
\frac{1}{\mathrm{Re}\,\mathrm{Pr}}
\frac{
T_{i+1,j}-2T_{ij}+T_{i-1,j}
}{\Delta x^2},
\end{aligned}
\label{eq:x_operator}
\end{equation}
and
\begin{equation}
\begin{aligned}
\left(\mathcal{Y}T\right)_{ij}
={}&
\frac{
v_{i,j+1}T_{i,j+1}
-v_{i,j-1}T_{i,j-1}
}{2\Delta y}
\\
&-
\frac{1}{\mathrm{Re}\,\mathrm{Pr}}
\frac{
T_{i,j+1}-2T_{ij}+T_{i,j-1}
}{\Delta y^2}.
\end{aligned}
\label{eq:y_operator}
\end{equation}
With these definitions, the non-source part of the spatial
operator satisfies
\begin{equation}
FI_{ij}
=
-\left(\mathcal{X}T\right)_{ij}
-\left(\mathcal{Y}T\right)_{ij}.
\label{eq:fi_operator_relation}
\end{equation}

Multiplying Equation~\eqref{eq:delta_form} by $\Delta t$
gives
\begin{equation}
\left(
I+\Delta t\,\mathcal{X}
+\Delta t\,\mathcal{Y}
\right)\delta T^{n+1}
=
\Delta t\,RHS^n,
\label{eq:implicit_delta_operator}
\end{equation}
where $I$ denotes the identity operator. This equation is
the operator form of the implicit Euler update. Directly
solving it requires a coupled two-dimensional system because
both directional operators act on the unknown temperature
update.

\subsubsection{Approximate Factorization}

The approximate factorization method replaces the combined
operator in Equation~\eqref{eq:implicit_delta_operator}
with a product of two directional operators:
\begin{equation}
\boxed{
\left(I+\Delta t\,\mathcal{X}\right)
\left(I+\Delta t\,\mathcal{Y}\right)
\delta T^{n+1}
=
\Delta t\,RHS^n.
}
\label{eq:approximate_factorization}
\end{equation}

Expanding the product on the left-hand side gives
\begin{equation}
\begin{aligned}
\Big[
I+\Delta t\,\mathcal{X}
+\Delta t\,\mathcal{Y}
+\Delta t^2\mathcal{X}\mathcal{Y}
\Big]\delta T^{n+1}
=
\Delta t\,RHS^n.
\end{aligned}
\label{eq:factorization_expanded}
\end{equation}

Comparing Equation~\eqref{eq:factorization_expanded} with
Equation~\eqref{eq:implicit_delta_operator} shows that
the factorized formulation introduces the additional
cross term
\begin{equation}
\Delta t^2\mathcal{X}\mathcal{Y}\delta T^{n+1}.
\label{eq:factorization_error_term}
\end{equation}
This term is proportional to $\Delta t^2$. The approximation
is therefore reasonable when the time step is sufficiently
small for the additional factorization term to remain
acceptable. However, the temporal accuracy of the underlying
implicit Euler method remains first-order.

The main advantage of this approximation is computational.
Instead of solving one coupled two-dimensional system, the
factorized equation can be solved through two successive
directional systems. This procedure is commonly implemented
using an alternating-direction implicit (ADI) strategy.

\subsubsection{Alternating-Direction Implicit Solution}

To solve Equation~\eqref{eq:approximate_factorization}
sequentially, an intermediate temperature update
$\widetilde{\delta T}^{\,n+1}$ is introduced as
\begin{equation}
\widetilde{\delta T}^{\,n+1}
=
\left(I+\Delta t\,\mathcal{Y}\right)
\delta T^{n+1}.
\label{eq:intermediate_update}
\end{equation}

Substituting this definition into
Equation~\eqref{eq:approximate_factorization} yields
the first directional system:
\begin{equation}
\boxed{
\left(I+\Delta t\,\mathcal{X}\right)
\widetilde{\delta T}^{\,n+1}
=
\Delta t\,RHS^n.
}
\label{eq:adi_step_x}
\end{equation}
After solving for the intermediate update, the second
directional system is solved:
\begin{equation}
\boxed{
\left(I+\Delta t\,\mathcal{Y}\right)
\delta T^{n+1}
=
\widetilde{\delta T}^{\,n+1}.
}
\label{eq:adi_step_y}
\end{equation}

The first step treats the $x$-direction operator implicitly,
while the second treats the $y$-direction operator
implicitly. Since each directional operator couples a
cell only to its immediate neighbors along one coordinate
direction, the resulting linear systems have tridiagonal
coefficient matrices along the respective grid lines.
These systems can be solved efficiently using a
tridiagonal matrix algorithm.

Finally, after obtaining the temperature update, the
temperature field is advanced according to
\begin{equation}
\boxed{
T_{ij}^{n+1}=T_{ij}^{n}+\delta T_{ij}^{n+1}.
}
\label{eq:temperature_update_final}
\end{equation}

The approximate factorization method thus reduces the
solution of a two-dimensional implicit problem to two
successive one-dimensional solves. When the objective is
to obtain a steady-state temperature field, this procedure
can be used for pseudo-time marching until the residual
falls below the prescribed convergence tolerance.

\paragraph{Treatment of Boundary Conditions}

The ghost-cell relations derived in
Section~\ref{subsec:ghost_cell_boundary_conditions}, including
their incremental forms, are incorporated into the directional
systems before they are solved. Since the unknown is the
temperature increment $\delta T^{n+1}$, the prescribed inlet
and wall temperatures lead to homogeneous incremental
boundary relations, whereas the zero-gradient outlet condition
requires the ghost-cell increment to equal the adjacent
physical-cell increment. Substituting these relations into
the directional operators eliminates ghost-cell increments
as independent unknowns and modifies the coefficients of
the corresponding tridiagonal systems. After solving the
two directional systems successively, the temperature field
is updated using
\begin{equation}
T_{ij}^{n+1}
=
T_{ij}^{n}+\delta T_{ij}^{n+1}.
\end{equation}

\subsection{Explicit Runge--Kutta Methods}
\label{subsec:runge_kutta}

The spatial discretization of the energy equation yields a system
of ordinary differential equations of the form
\begin{equation}
\frac{dT_{ij}}{dt}=RHS_{ij}(T)=FI_{ij}(T)+S_{ij},
\label{eq:rk_semidiscrete}
\end{equation}
where $FI_{ij}(T)$ represents the discretized advective and
diffusive contributions, and $S_{ij}$ is the viscous-dissipation
source term. For convenience, the complete right-hand side is
denoted by the operator $\mathcal{F}$:
\begin{equation}
\mathcal{F}_{ij}(T)=FI_{ij}(T)+S_{ij}.
\label{eq:rk_rhs_operator}
\end{equation}

The Runge--Kutta methods advance the solution explicitly by
evaluating this operator at several intermediate stages within
each time step. Unlike the approximate factorization method,
they do not require the solution of implicit linear systems.
The second- and fourth-order methods considered in this project
are described below.

\subsubsection{Second-Order Runge--Kutta Method}

The second-order Runge--Kutta method used in this project is
the explicit midpoint method. It evaluates the right-hand side
at the beginning of the time step and at an intermediate
state located halfway through the step.

First, the initial slope is evaluated using the temperature
field at time level $n$:
\begin{equation}
k_{1,ij}=\mathcal{F}_{ij}(T^n).
\label{eq:rk2_k1}
\end{equation}

An intermediate temperature field is then constructed using
a half time step:
\begin{equation}
T_{ij}^{(1)}
=
T_{ij}^{n}+\frac{\Delta t}{2}k_{1,ij}.
\label{eq:rk2_intermediate}
\end{equation}

The right-hand side is evaluated again at this intermediate
state:
\begin{equation}
k_{2,ij}=\mathcal{F}_{ij}\left(T^{(1)}\right).
\label{eq:rk2_k2}
\end{equation}

The solution at the next time level is obtained using the
intermediate slope:
\begin{equation}
\boxed{
T_{ij}^{n+1}
=
T_{ij}^{n}+\Delta t\,k_{2,ij}.
}
\label{eq:rk2_update}
\end{equation}

The explicit midpoint method is second-order accurate in
time. Its use of an intermediate evaluation provides a more
accurate approximation to the evolution of the solution over
a time step than the first-order explicit Euler method.

\subsubsection{Fourth-Order Runge--Kutta Method}

The classical fourth-order Runge--Kutta method evaluates
the right-hand side four times during each time step. These
evaluations provide estimates of the solution's slope at
the beginning, at two intermediate states, and at the end
of the time step.

The first stage evaluates the right-hand side at the
current temperature field:
\begin{equation}
k_{1,ij}=\mathcal{F}_{ij}(T^n).
\label{eq:rk4_k1}
\end{equation}

The second stage evaluates it at an intermediate state
predicted using the first slope:
\begin{equation}
k_{2,ij}
=
\mathcal{F}_{ij}
\left(
T^n+\frac{\Delta t}{2}k_1
\right).
\label{eq:rk4_k2}
\end{equation}

The third stage uses the second slope to construct another
intermediate state:
\begin{equation}
k_{3,ij}
=
\mathcal{F}_{ij}
\left(
T^n+\frac{\Delta t}{2}k_2
\right).
\label{eq:rk4_k3}
\end{equation}

Finally, the fourth stage evaluates the right-hand side
at a state predicted over the full time step:
\begin{equation}
k_{4,ij}
=
\mathcal{F}_{ij}
\left(
T^n+\Delta t\,k_3
\right).
\label{eq:rk4_k4}
\end{equation}

The four stage evaluations are combined to obtain the
updated temperature field:
\begin{equation}
\boxed{
T_{ij}^{n+1}
=
T_{ij}^{n}
+\frac{\Delta t}{6}
\left(
k_{1,ij}
+2k_{2,ij}
+2k_{3,ij}
+k_{4,ij}
\right).
}
\label{eq:rk4_update}
\end{equation}

The classical Runge--Kutta method is fourth-order accurate
in time for sufficiently smooth solutions. Although it
requires four evaluations of the spatial operator per time
step, its higher temporal accuracy can be advantageous
when the time step is sufficiently small and stability
requirements are satisfied.

\subsection{Direct Sparse Solution}
\label{subsec:direct_solution}

\subsubsection{Linearization of the Spatial Operator}

The spatially discretized energy equation can be expressed as
a system of ordinary differential equations:
\begin{equation}
\frac{d\mathbf{T}}{dt}
=
\mathcal{L}(\mathbf{T}),
\label{eq:direct_ode}
\end{equation}
where $\mathbf{T}$ is the vector of cell-centered temperatures
and $\mathcal{L}$ denotes the spatial operator, including
advection, diffusion, and viscous dissipation.

For the direct solver, the spatial operator is written in the
form
\begin{align}
\mathcal{L}_{ij}(T)
={}&
-c\frac{
u_{i+1,j}T_{i+1,j}
-u_{i-1,j}T_{i-1,j}
}{2\Delta x}
\nonumber\\
&-
c\frac{
v_{i,j+1}T_{i,j+1}
-v_{i,j-1}T_{i,j-1}
}{2\Delta y}
\nonumber\\
&+
\frac{1}{\mathrm{Pe}}
\frac{
T_{i+1,j}-2T_{i,j}+T_{i-1,j}
}{\Delta x^2}
\nonumber\\
&+
\frac{1}{\mathrm{Pe}}
\frac{
T_{i,j+1}-2T_{i,j}+T_{i,j-1}
}{\Delta y^2}
+S_{ij},
\label{eq:direct_spatial_operator}
\end{align}
where $\mathrm{Pe}$ is the Peclet number, $c$ is the convection
coefficient, and $S_{ij}$ is the viscous-dissipation source
term. The velocity field and source term are prescribed and
independent of temperature. Consequently, the discretized
spatial operator is linear in the unknown temperatures,
provided that the boundary conditions are linear.

\subsubsection{Sparse Matrix Formulation}

Collecting the temperatures of the $n_xn_y$ physical cells
into a vector gives
\begin{equation}
\boxed{
\frac{d\mathbf{T}}{dt}
=
\mathbf{A}\mathbf{T}+\mathbf{b},
}
\label{eq:direct_semidiscrete}
\end{equation}
where $\mathbf{A}$ is the sparse matrix representing the
linear spatial operator, and $\mathbf{b}$ contains the
temperature-independent contributions from the boundary
conditions and viscous dissipation.

The matrix is sparse because each discretized equation
couples a cell primarily to its immediate neighbors. The
physical-cell temperature vector is constructed by flattening
the two-dimensional temperature array in C order, with the
wall-normal index varying fastest. The corresponding
zero-based array indices are mapped to the vector index
\begin{equation}
k(i,j)=i n_y+j,
\qquad
i=0,\ldots,n_x-1,\qquad
j=0,\ldots,n_y-1.
\label{eq:direct_vector_index}
\end{equation}

The forcing vector can be obtained by evaluating the spatial
operator with all physical-cell temperatures set to zero,
while retaining the prescribed boundary conditions:
\begin{equation}
\boxed{
\mathbf{b}=\mathcal{L}(\mathbf{0}).
}
\label{eq:direct_forcing_vector}
\end{equation}
The ghost-cell relations derived in
Section~\ref{subsec:ghost_cell_boundary_conditions} are substituted
into the spatially discretized equations before assembling the
coefficient matrix. This eliminates ghost-cell temperatures as
independent unknowns. Contributions proportional to physical-cell
temperatures are incorporated into the matrix $\mathbf{A}$,
while temperature-independent contributions from prescribed
boundary values are included in the forcing vector $\mathbf{b}$.
Consequently, the resulting linear system contains only the
temperatures of the physical cells.

\subsubsection{Implicit Euler Time Integration}

The semi-discrete system is advanced in pseudo-time using
the implicit Euler method:
\begin{equation}
\frac{\mathbf{T}^{n+1}-\mathbf{T}^{n}}{\Delta t}
=
\mathbf{A}\mathbf{T}^{n+1}+\mathbf{b},
\label{eq:direct_implicit_euler}
\end{equation}
where $\Delta t$ denotes the pseudo-time step. Rearranging
the equation gives the linear system
\begin{equation}
\boxed{
\left(\mathbf{I}-\Delta t\,\mathbf{A}\right)
\mathbf{T}^{n+1}
=
\mathbf{T}^{n}+\Delta t\,\mathbf{b}.
}
\label{eq:direct_linear_system}
\end{equation}

Defining the coefficient matrix as
\begin{equation}
\mathbf{M}=\mathbf{I}-\Delta t\,\mathbf{A},
\label{eq:direct_matrix_definition}
\end{equation}
the update becomes
\begin{equation}
\mathbf{M}\mathbf{T}^{n+1}
=
\mathbf{T}^{n}+\Delta t\,\mathbf{b}.
\label{eq:direct_matrix_update}
\end{equation}

When the grid, pseudo-time step, velocity field, model
parameters, and boundary conditions remain fixed, the matrix
$\mathbf{M}$ does not change between iterations. Its sparse
LU factorization can therefore be computed once and reused
for subsequent time steps:
\begin{equation}
\mathbf{M}=\mathbf{L}_{\mathrm{LU}}\mathbf{U}_{\mathrm{LU}}.
\label{eq:direct_lu_factorization}
\end{equation}
At each iteration, the right-hand side is first solved using
the lower-triangular factor:
\begin{equation}
\mathbf{L}_{\mathrm{LU}}\mathbf{y}
=
\mathbf{T}^{n}+\Delta t\,\mathbf{b},
\label{eq:direct_forward_substitution}
\end{equation}
followed by the upper-triangular solve
\begin{equation}
\mathbf{U}_{\mathrm{LU}}\mathbf{T}^{n+1}
=
\mathbf{y}.
\label{eq:direct_backward_substitution}
\end{equation}

This procedure uses a direct linear-system solver at each
time step rather than an iterative method to solve the
coupled linear equations. Reusing the LU factorization avoids
repeating the relatively expensive matrix-factorization
stage during the pseudo-time-marching process.

In summary, the direct solver combines a sparse matrix
representation of the spatially discretized energy equation
with implicit Euler pseudo-time marching and a reusable sparse
LU factorization. This approach provides a reference numerical
solution against which the approximate factorization and
explicit Runge--Kutta methods can be compared.

\subsection{Convergence Assessment}
\label{sec:convergence_assessment}

The objective of the numerical time-marching methods is to obtain
a steady-state temperature distribution. At steady state, the
temporal derivative vanishes, and the spatially discretized
right-hand side approaches zero:
\begin{equation}
RHS_{ij}(T)=FI_{ij}(T)+S_{ij}=0.
\label{eq:steady_state_equation}
\end{equation}

Convergence is assessed using the root-mean-square (RMS)
residual of the spatially discretized energy equation over
the physical cells:
\begin{equation}
\boxed{
R_{\mathrm{RMS}}
=
\sqrt{
\frac{1}{n_xn_y}
\sum_{i=1}^{n_x}
\sum_{j=1}^{n_y}
\left[RHS_{ij}(T)\right]^2
}.
}
\label{eq:rms_residual}
\end{equation}

This residual provides a common convergence measure for all
numerical solution methods considered in this project, since
it evaluates how closely the current temperature field
satisfies the same discrete steady-state energy equation,
independently of the temporal integration scheme.

The calculation is terminated when the RMS residual falls
below a prescribed tolerance or the maximum number of
iterations is reached. The residual is therefore used to
assess convergence toward a steady-state numerical solution.
However, a small residual does not by itself quantify the
spatial discretization error or establish grid independence;
these require separate verification studies.